\documentclass{article}

% if you need to pass options to natbib, use, e.g.:
%     \PassOptionsToPackage{numbers, compress}{natbib}
% before loading neurips_2026

% The authors should use one of these tracks.
% Before accepting by the NeurIPS conference, select one of the options below.
% 0. "default" for submission
\usepackage{neurips_2026}
% the "default" option is equal to the "main" option, which is used for the Main Track with double-blind reviewing.
% 1. "main" option is used for the Main Track
%  \usepackage[main]{neurips_2026}
% 2. "position" option is used for the Position Paper Track
%  \usepackage[position]{neurips_2026}
% 3. "eandd" option is used for the Evaluations & Datasets Track
 % \usepackage[eandd]{neurips_2026}
 % if you want to opt-in for a de-anomymized submission (not recommended, but OK):
 %\usepackage[eandd, nonanonymous]{neurips_2026}
% 4. "creativeai" option is used for the Creative AI Track
%  \usepackage[creativeai]{neurips_2026}
% 5. "sglblindworkshop" option is used for the Workshop with single-blind reviewing
 % \usepackage[sglblindworkshop]{neurips_2026}
% 6. "dblblindworkshop" option is used for the Workshop with double-blind reviewing
%  \usepackage[dblblindworkshop]{neurips_2026}

% After being accepted, the authors should add "final" behind the track to compile a camera-ready version.
% 1. Main Track
 % \usepackage[main, final]{neurips_2026}
% 2. Position Paper Track
%  \usepackage[position, final]{neurips_2026}
% 3. Evaluations & Datasets Track
 % \usepackage[eandd, final]{neurips_2026}
% 4. Creative AI Track
%  \usepackage[creativeai, final]{neurips_2026}
% 5. Workshop with single-blind reviewing
%  \usepackage[sglblindworkshop, final]{neurips_2026}
% 6. Workshop with double-blind reviewing
%  \usepackage[dblblindworkshop, final]{neurips_2026}
% Note. For the workshop paper template, both \title{} and \workshoptitle{} are required, with the former indicating the paper title shown in the title and the latter indicating the workshop title displayed in the footnote.
% For workshops (5., 6.), the authors should add the name of the workshop, "\workshoptitle" command is used to set the workshop title.
% \workshoptitle{WORKSHOP TITLE}

% "preprint" option is used for arXiv or other preprint submissions
 % \usepackage[preprint]{neurips_2026}

% to avoid loading the natbib package, add option nonatbib:
%    \usepackage[nonatbib]{neurips_2026}

\usepackage[utf8]{inputenc} % allow utf-8 input
\usepackage[T1]{fontenc}    % use 8-bit T1 fonts
\usepackage{hyperref}       % hyperlinks
\usepackage{url}            % simple URL typesetting
\usepackage{booktabs}       % professional-quality tables
\usepackage{amsfonts}       % blackboard math symbols
\usepackage{nicefrac}       % compact symbols for 1/2, etc.
\usepackage{microtype}      % microtypography
\usepackage{xcolor}         % colors

% Note. For the workshop paper template, both \title{} and \workshoptitle{} are required, with the former indicating the paper title shown in the title and the latter indicating the workshop title displayed in the footnote. 
\title{Toolflix: Feedback-Driven Tool Recommendation for Agentic Systems}

\author{%
  % TODO: Add authors
}

\begin{document}

\maketitle

\begin{abstract}
% TODO: Write after the paper is done. Should cover: problem (MCP tool explosion,
% description-reality gap), approach (two-stage retrieve-then-rerank with two-head
% wide \& deep model), key result (learning curve showing quality head lifts
% success@k over retriever-only baseline as feedback accumulates).
\end{abstract}


%% ============================================================
\section{Introduction}
\label{sec:intro}
%% ============================================================

% P1: Context + problem. One idea per sentence.
Large language models are increasingly augmented with external tools to act as autonomous agents~\citep{schick2024toolformer, yao2023react, shen2024hugginggpt}. Most production agents today operate over curated toolsets where a human has already vetted each tool. But the ecosystem is shifting toward open marketplaces: standardized protocols like MCP~\citep{anthropic2025mcp} now expose thousands of community-contributed endpoints, and retrieval-augmented approaches already struggle at this scale---tool selection accuracy collapses from 43\% to under 14\% as registries grow~\citep{sun2025ragmcp}. The core difficulty is not finding tools that \emph{match} a query, but finding tools that \emph{work}. In open registries, multiple endpoints claim to serve the same function with near-identical descriptions, yet they differ sharply in reliability. We call this \emph{underdifferentiation}: tools that are semantically indistinguishable but qualitatively distinct.

% P2: Why existing work fails. Argument first, citations as evidence.
No existing approach can resolve underdifferentiation, because all operate on static, description-level information. In curated settings this was sufficient---a human had already filtered for quality. In open marketplaces, quality variance is the default, and descriptions carry no signal about it. Retrieval methods rank by semantic similarity and are blind when descriptions converge~\citep{shi2024toolret}. Active discovery improves query formulation but still routes by similarity~\citep{fei2025mcpzero}. Recommendation approaches use static benchmark labels with no runtime signal~\citep{shi2026agentselect}. Deployed systems like Anthropic's tool search use BM25 over metadata~\citep{anthropic2025toolsearch}. The gap is consistent: none of these systems learn from whether a tool actually succeeded.

% P3: Our approach. High-level mechanism only; details go in Section 4.
We propose \textsc{Toolflix}, a recommender that learns which tools work from execution feedback. \textsc{Toolflix} places a learned reranker after a standard semantic retriever. The reranker has two heads: a \emph{relevance head} that learns query-tool affinity (e.g., fetch endpoints work for search tasks, but specific ones fail on certain domains), and a \emph{quality head} that learns per-tool reliability from accumulated success rates---the equivalent of a star rating that emerges from usage, not descriptions. A blending parameter $\alpha$ shifts from relevance-dominant at cold start to quality-dominant as feedback accumulates. After each task, an LLM judge provides a binary success/failure signal; training uses pairwise margin loss over these outcomes and requires no human labels.

% P4: Key evidence preview. TODO: fill with final numbers.
% Testbed: 59 real + 104 synthetic MCP endpoints with controlled failure modes.
% Evaluation axes: (1) learning curve, (2) held-out comparison,
% (3) cold-start/OOD, (4) quality drift, (5) transfer to TOOLRET.
% Show underdifferentiation score correlates with reranker improvement.

% P5: Contributions. One sentence each.
Our contributions:
\begin{itemize}
    \item We formalize \emph{underdifferentiation}---high description similarity coupled with high variance in execution quality---and show that the reranker's gain over retrieval correlates with this measurable score.
    \item We propose a two-head reranker that separates learned relevance from learned quality, with an adaptive blending schedule for cold start.
    \item We release a synthetic tool benchmark with controlled quality variance: paraphrased descriptions and injected failure modes that are invisible to retrievers, isolating the value of feedback.
    \item We show feedback-driven reranking improves selection over retrieval baselines, with the largest gains in high-underdifferentiation categories.
\end{itemize}


%% ============================================================
\section{Related work}
\label{sec:related}
%% ============================================================

\subsection{Tool-augmented LLMs}

Tool integration has progressed from hard-coded access~\citep{yao2023react} to learned invocation~\citep{schick2024toolformer, patil2023gorilla} to context-based function calling~\citep{shen2024hugginggpt}. Benchmarks like ToolBench~\citep{qin2024toolbench} and API-Bank~\citep{li2023apibank} evaluate tool \emph{usage}---can the agent call a tool correctly?---but pre-annotate 10--20 candidate tools per task, bypassing the selection problem entirely. We focus on the upstream question: which tool to recommend from a large, unvetted pool.

\subsection{Tool retrieval and selection}

Several lines of work address tool selection at scale, each improving a different part of the pipeline but none incorporating execution feedback.

\paragraph{Better retrieval.} \citet{shi2024toolret} show that standard IR models perform poorly on tool corpora due to low lexical overlap and domain shift, and that tool-specific training data substantially improves both retrieval and downstream task pass rates. RAG-MCP~\citep{sun2025ragmcp} applies retrieval-augmented pre-filtering to MCP tools but degrades sharply beyond ${\sim}100$ tools. Tool-Embed~\citep{toolembed2026} uses LLM-driven document expansion to achieve state-of-the-art on TOOLRET, observing that identical functions admit seven or more distinct phrasings---independent evidence for underdifferentiation at the description level.

\paragraph{Better reranking.} ToolRerank~\citep{zheng2024toolrerank} introduces hierarchy-aware reranking using structural priors. ScaleMCP~\citep{lumer2025scalemcp} finds that LLM-based reranking over 5,000 MCP tools is effective but computationally expensive. Our two-head reranker is lightweight by comparison, trained on implicit feedback rather than requiring LLM inference at ranking time.

\paragraph{Better queries.} MCP-Zero~\citep{fei2025mcpzero} has the agent generate structured tool requests that are better aligned with tool documentation than raw user queries. Anthropic's tool search~\citep{anthropic2025toolsearch} uses BM25 or regex matching with deferred loading, reducing tokens by 85\% and raising accuracy from 49\% to 74\%. Both solve \emph{discovery}---which tools to surface---but not \emph{quality selection}: when two tools have near-identical descriptions, similarity-based methods rank them identically.

\paragraph{Beyond tool endpoints.} In the adjacent domain of agent skill files---long-form behavioral instructions rather than callable endpoints---SkillRouter~\citep{zheng2026skillrouter} finds that metadata alone is insufficient to route among 80K community skills with pervasive functional overlap, with a 29--44pp drop when the full implementation body is removed. Though skills differ fundamentally from tool endpoints (they modify agent behavior rather than returning structured outputs, making execution-based evaluation difficult), the finding reinforces our thesis: description-level information is insufficient for selection whenever a registry contains many semantically similar entries.

The consistent gap across all of the above: a tool that matches perfectly but fails 80\% of the time will be retrieved indefinitely.

\subsection{Agent and tool recommendation}

\citet{shi2026agentselect} frame agent selection as recommendation, ranking full (model + toolset) configurations from static benchmark labels. Their finding that content-aware matching outperforms collaborative filtering in the long tail parallels our underdifferentiation setting. They identify ``selective end-to-end execution feedback'' as future work; our approach implements this directly.

\subsection{Learning from implicit feedback}

Our reranker combines ideas from implicit-feedback ranking and the cold-start problem. BPR~\citep{rendle2009bpr} established pairwise optimization from implicit signals; the wide-and-deep architecture~\citep{cheng2016wide} showed that combining memorization (per-item statistics) with generalization (learned embeddings) outperforms either alone. We adopt this decomposition: wide features memorize per-tool success rates, deep features generalize query-tool affinity to unseen queries.

New tools present a cold-start problem: no interaction history means no quality signal. Contextual bandits~\citep{li2010contextual} and Thompson sampling~\citep{chapelle2011thompson} address this via uncertainty-driven exploration. Our $\alpha$-blending schedule serves an analogous role---falling back to the relevance head when quality estimates are unreliable---complemented by a UCB-style exploration bonus for underexplored tools. This avoids full posterior inference while providing a principled exploration schedule.


%% ============================================================
\section{Problem setting}
\label{sec:problem}
%% ============================================================

% Formalize the setting:
%   - Pool of N tool endpoints, each with description and JSON schema
%   - Agent receives user task, decomposes into query
%   - System returns top-k tools; agent selects one and executes
%   - Agent provides feedback: relevance (binary), success (binary), score (1-5)
%
% Define the ranking objective: maximize success@k over a stream of tasks.
% Distinguish from the agent's tool-calling ability (orthogonal).


%% ============================================================
\section{Method}
\label{sec:method}
%% ============================================================

Our system operates in two stages: a \emph{retrieval} stage that maximizes recall over the full tool pool, and a \emph{reranking} stage that maximizes precision using learned feedback signals. An agentic loop connects these stages, executing tools and collecting the feedback that drives online learning.

\subsection{Stage 1: Hybrid retrieval}

Given a user task, an LLM agent first decomposes it into a short tool-capability query $q$ (e.g., ``extract text from a local PDF file''). We retrieve the top-$K$ candidate tool endpoints by hybrid scoring:
%
\begin{equation}
  \text{score}(q, t) = 0.7 \cdot \cos\bigl(\phi(q),\, \phi(d_t)\bigr) + 0.3 \cdot \text{kw}(q, d_t)
\end{equation}
%
where $\phi$ denotes the sentence-transformer encoder (all-MiniLM-L6-v2, 384-d), $d_t$ is the tool description text, and $\text{kw}(\cdot)$ is a token-overlap score that catches lexical matches the embedding model misses. We set $K{=}100$ to ensure high recall: with 163 endpoints, this covers 61\% of the pool and ensures that even tools with atypical descriptions (e.g., a search tool whose description emphasizes non-English engines) are surfaced for the reranker.

\subsection{Stage 2: Wide \& deep reranker}

The reranker scores each of the $K$ candidates and returns the top 5 to the agent. It follows a two-head wide-and-deep architecture with 68K parameters, inspired by~\citet{cheng2016wide} but adapted for pairwise ranking from implicit feedback.

\paragraph{Shared deep trunk.} For a query $q$ and candidate tool $t$ with description embedding $\phi(d_t)$, the trunk computes interaction features between learned projections:
%
\begin{align}
  \mathbf{q} &= W_q\, \phi(q) \in \mathbb{R}^{64}, \quad \mathbf{d} = W_d\, \phi(d_t) \in \mathbb{R}^{64} \\
  \mathbf{e}_t &\in \mathbb{R}^{32} \quad \text{(learned endpoint embedding, initialized from } \phi(d_t)\text{)}
\end{align}
%
The trunk concatenates four signals---cosine similarity $\cos(\mathbf{q}, \mathbf{d})$, element-wise product $\mathbf{q} \odot \mathbf{d}$, difference $\mathbf{q} - \mathbf{d}$, and $\mathbf{e}_t$---and maps them through a two-layer MLP ($161 \to 64 \to 32$) with ReLU activations and dropout (0.2), producing a deep representation $\mathbf{h} \in \mathbb{R}^{32}$.

The learned endpoint embedding $\mathbf{e}_t$ is the key component. Initialized from the description, it diverges through training to encode what the tool \emph{actually does} rather than what its description claims. The projection matrices $W_q$ and $W_d$ are shared across all tools, learning \emph{how to read descriptions skeptically}---amplifying dimensions that predict execution success and suppressing those that don't.

\paragraph{Relevance head (personalized recommendation).} A two-layer MLP ($32 \to 16 \to 1$) that reads \emph{only} the deep representation $\mathbf{h}$:
%
\begin{equation}
  r(q, t) = \text{MLP}_{\text{rel}}(\mathbf{h})
\end{equation}
%
With no access to per-tool execution statistics, this head is forced to learn query-tool affinity entirely through the embeddings. It captures \emph{query-conditioned} quality: a tool may have high overall success but fail for a specific query type. By analogy with collaborative filtering, this head provides personalized recommendations---``given \emph{your query}, this tool works best''---even for tools with limited execution history.

\paragraph{Quality head (star rating).} A two-layer MLP ($38 \to 16 \to 1$) that reads the deep representation \emph{concatenated with} six wide features:
%
\begin{equation}
  s(q, t) = \text{MLP}_{\text{qual}}\bigl([\mathbf{w}_t \,\|\, \mathbf{h}]\bigr)
\end{equation}
%
The wide feature vector $\mathbf{w}_t \in \mathbb{R}^6$ contains: normalized usage count, success rate, relevance rate, query-category co-occurrence, category match, and retriever similarity. These are aggregate execution statistics updated after every task at $O(1)$ cost. The quality head memorizes per-tool reliability---``this tool succeeds 90\% of the time''---providing the global star rating.

\paragraph{Score combination.} The final ranking score blends both heads with an adaptive weight $\alpha$ that decays as feedback accumulates:
%
\begin{equation}
  \text{score}(q, t) = \alpha \cdot r(q, t) + (1 - \alpha) \cdot s(q, t), \quad \alpha = 0.35 + 0.45 \cdot e^{-n/80}
\end{equation}
%
where $n$ is the total number of feedback signals received. At cold start ($n{=}0$), $\alpha \approx 0.80$: the system trusts the relevance head's embedding-based generalization. As feedback accumulates, $\alpha$ converges to $0.35$, shifting weight toward the quality head's execution statistics while retaining the relevance head's query-specific signal. We found $\alpha{=}0.35$ optimal at convergence; pure quality ($\alpha{=}0$) underperforms because it ignores query-tool affinity, while equal weighting ($\alpha{=}0.5$) gives too much influence to the noisier relevance signal.

An exploration bonus $\beta \cdot \sqrt{\log(n+1) / (n_t + 1)}$ is added for underexplored tools, where $\beta$ decays exponentially to zero as feedback accumulates.

\subsection{Training}

Both heads train on the same pairwise success signal: when a tool succeeds, it should rank above the other candidates; when it fails, others should rank above it. For each feedback record with candidates $\mathcal{C}$, selected tool $t^*$, and binary outcome $y$, we generate pairs:
%
\begin{equation}
  (t^+, t^-) = \begin{cases} (t^*, t_j) & \text{if } y = 1 \\ (t_j, t^*) & \text{if } y = 0 \end{cases}, \quad \forall\, t_j \in \mathcal{C} \setminus \{t^*\}
\end{equation}
%
Training uses weighted pairwise margin loss:
%
\begin{equation}
  \mathcal{L} = \frac{1}{|\mathcal{P}|} \sum_{(t^+, t^-) \in \mathcal{P}} w_{ij} \cdot \max\bigl(0,\; m - f(q, t^+) + f(q, t^-)\bigr)
\end{equation}
%
where $m{=}0.5$ is the margin and $w_{ij}{=}3$ for hard negatives (same-server pairs or both with similarity $> 0.35$), $w_{ij}{=}1$ otherwise. The key difference between heads: the relevance head sees $\mathbf{w}_t = \mathbf{0}$ (zeroed wide features), forcing it to learn through embeddings alone; the quality head sees real wide features.

Training proceeds in three phases per batch: (1) relevance head only (40\% of epochs), updating $W_q$, $W_d$, $\mathbf{e}_t$, trunk, and relevance MLP; (2) quality head only (40\%), updating trunk and quality MLP; (3) joint fine-tuning (20\%) at reduced learning rate ($0.3\times$). The learning rate adapts to dataset size: $\text{lr} = \min(5{\times}10^{-3},\; 5 / (n + 200))$.

\paragraph{Online learning.} Wide features update after every task. The deep model retrains on a replay buffer every 50 examples, with epochs scaled by buffer size ($\max(10, \min(15, 20 - n/200))$). This provides immediate incorporation of new feedback (via wide features) with periodic refinement of the deep representation.

\subsection{Agentic loop}

Each task requires three LLM calls:
\begin{enumerate}[leftmargin=*]
  \item \textbf{Decompose.} The agent receives the user task and outputs a short capability query for the retriever.
  \item \textbf{Select \& call.} The top-5 reranked tools are presented as native function-calling tools with match percentages derived from reranker scores ($30$--$99\%$ range). The agent selects one and fills its arguments; the system executes the tool via JSON-RPC over stdio.
  \item \textbf{Rate.} The agent evaluates the tool result, producing structured feedback: relevance (binary), success (binary), quality score (1--5), and reasoning.
\end{enumerate}
%
The feedback record---including the full candidate list, selected tool, arguments, result preview, and rating---is appended to the replay buffer and used for online learning.


%% ============================================================
\section{Experimental setup}
\label{sec:experiments}
%% ============================================================

% Tool pool: 163 endpoints (59 real across 16 MCP servers in 4 categories
% + 104 synthetic variants with paraphrased descriptions and injected
% failure modes). Describe the category-level quality variance that
% motivates the work (e.g., PDF tools range from 30-90\% success).
% Compute underdifferentiation score per category: avg pairwise cosine
% similarity of descriptions * variance in success rates. Show this
% correlates with reranker improvement.
%
% Task dataset: 1000 LLM-generated tasks with concrete artifacts.
% 200 held-out test tasks.
%
% Baselines: (1) BM25, (2) embedding retriever (cosine similarity),
% (3) success-rate prior (wide-only proxy), (4) random,
% (5) ToolRerank (hierarchy-aware, no feedback),
% (6) SkillRouter-style retrieve-and-rerank (description-based).
%
% Ablations: relevance head only, quality head only, combined,
% alpha fixed vs. adaptive schedule.
%
% Transfer experiment: train on our testbed, evaluate on TOOLRET
% categories with high description overlap. Shows generalization
% beyond custom testbed. Addresses scale gap vs. TOOLRET (43K tools),
% MCP-Zero (2,797 tools), AgentSelect (107K agents).
%
% Metrics: AUC, pairwise accuracy, success@1, success@5, avg rank of
% succeeded tools. Per-category breakdown. Underdifferentiation score
% correlation.


%% ============================================================
\section{Results}
\label{sec:results}
%% ============================================================

% Main result: learning curve (Figure 1). X-axis = feedback count,
% Y-axis = success@k. Show retriever (flat), relevance head (slight lift),
% quality head (climbing), combined (best). The gap between relevance-only
% and combined is the value of execution feedback. Cold start to convergence.
%
% Table 1: final metrics for all baselines and ablations.
%
% Table 2: per-category breakdown showing where feedback helps most
% (categories with high quality variance, e.g., search with DuckDuckGo
% at 37\% vs Brave at 80\%).
%
% Analysis of learned embeddings: show examples where endpoint embeddings
% diverged from descriptions (description-reality gap captured).


%% ============================================================
\section{Discussion}
\label{sec:discussion}
%% ============================================================

% What the two-head decomposition reveals: relevance head alone performs
% near retriever level (semantic features are necessary but not sufficient),
% quality head provides the marginal lift (feedback is the differentiator).
%
% Limitations: small endpoint pool (52), synthetic tasks, single agent model.
% Feedback quality depends on the rating agent. Cold start requires
% exploration budget.
%
% Future work: scale to 1000+ endpoints, multi-agent feedback aggregation,
% listwise loss (LambdaRank), serve as live API.


%% ============================================================
\section{Conclusion}
\label{sec:conclusion}
%% ============================================================

% Restate the core finding: execution feedback, decomposed into a relevance
% head and a quality head, enables a tool recommender that improves from
% cold start and outperforms static semantic retrieval. The quality head
% is the mechanism that captures what descriptions cannot.


\bibliography{references}
\bibliographystyle{plainnat}


%%%%%%%%%%%%%%%%%%%%%%%%%%%%%%%%%%%%%%%%%%%%%%%%%%%%%%%%%%%%

\appendix

\section{Tool endpoint details}
% Full list of 13 MCP servers and 52 endpoints with categories.

\section{Task generation details}
% Templates, artifact distributions, category breakdown.

\section{Training hyperparameters}
% Learning rates, margins, epoch schedules, batch sizes.

%%%%%%%%%%%%%%%%%%%%%%%%%%%%%%%%%%%%%%%%%%%%%%%%%%%%%%%%%%%%

\newpage
\input{checklist.tex}


\end{document}